% Options for packages loaded elsewhere
\PassOptionsToPackage{unicode}{hyperref}
\PassOptionsToPackage{hyphens}{url}
\documentclass[
]{article}
\usepackage{xcolor}
\usepackage[margin=1in]{geometry}
\usepackage{amsmath,amssymb}
\setcounter{secnumdepth}{-\maxdimen} % remove section numbering
\usepackage{iftex}
\ifPDFTeX
  \usepackage[T1]{fontenc}
  \usepackage[utf8]{inputenc}
  \usepackage{textcomp} % provide euro and other symbols
\else % if luatex or xetex
  \usepackage{unicode-math} % this also loads fontspec
  \defaultfontfeatures{Scale=MatchLowercase}
  \defaultfontfeatures[\rmfamily]{Ligatures=TeX,Scale=1}
\fi
\usepackage{lmodern}
\ifPDFTeX\else
  % xetex/luatex font selection
\fi
% Use upquote if available, for straight quotes in verbatim environments
\IfFileExists{upquote.sty}{\usepackage{upquote}}{}
\IfFileExists{microtype.sty}{% use microtype if available
  \usepackage[]{microtype}
  \UseMicrotypeSet[protrusion]{basicmath} % disable protrusion for tt fonts
}{}
\makeatletter
\@ifundefined{KOMAClassName}{% if non-KOMA class
  \IfFileExists{parskip.sty}{%
    \usepackage{parskip}
  }{% else
    \setlength{\parindent}{0pt}
    \setlength{\parskip}{6pt plus 2pt minus 1pt}}
}{% if KOMA class
  \KOMAoptions{parskip=half}}
\makeatother
\usepackage{color}
\usepackage{fancyvrb}
\newcommand{\VerbBar}{|}
\newcommand{\VERB}{\Verb[commandchars=\\\{\}]}
\DefineVerbatimEnvironment{Highlighting}{Verbatim}{commandchars=\\\{\}}
% Add ',fontsize=\small' for more characters per line
\usepackage{framed}
\definecolor{shadecolor}{RGB}{248,248,248}
\newenvironment{Shaded}{\begin{snugshade}}{\end{snugshade}}
\newcommand{\AlertTok}[1]{\textcolor[rgb]{0.94,0.16,0.16}{#1}}
\newcommand{\AnnotationTok}[1]{\textcolor[rgb]{0.56,0.35,0.01}{\textbf{\textit{#1}}}}
\newcommand{\AttributeTok}[1]{\textcolor[rgb]{0.13,0.29,0.53}{#1}}
\newcommand{\BaseNTok}[1]{\textcolor[rgb]{0.00,0.00,0.81}{#1}}
\newcommand{\BuiltInTok}[1]{#1}
\newcommand{\CharTok}[1]{\textcolor[rgb]{0.31,0.60,0.02}{#1}}
\newcommand{\CommentTok}[1]{\textcolor[rgb]{0.56,0.35,0.01}{\textit{#1}}}
\newcommand{\CommentVarTok}[1]{\textcolor[rgb]{0.56,0.35,0.01}{\textbf{\textit{#1}}}}
\newcommand{\ConstantTok}[1]{\textcolor[rgb]{0.56,0.35,0.01}{#1}}
\newcommand{\ControlFlowTok}[1]{\textcolor[rgb]{0.13,0.29,0.53}{\textbf{#1}}}
\newcommand{\DataTypeTok}[1]{\textcolor[rgb]{0.13,0.29,0.53}{#1}}
\newcommand{\DecValTok}[1]{\textcolor[rgb]{0.00,0.00,0.81}{#1}}
\newcommand{\DocumentationTok}[1]{\textcolor[rgb]{0.56,0.35,0.01}{\textbf{\textit{#1}}}}
\newcommand{\ErrorTok}[1]{\textcolor[rgb]{0.64,0.00,0.00}{\textbf{#1}}}
\newcommand{\ExtensionTok}[1]{#1}
\newcommand{\FloatTok}[1]{\textcolor[rgb]{0.00,0.00,0.81}{#1}}
\newcommand{\FunctionTok}[1]{\textcolor[rgb]{0.13,0.29,0.53}{\textbf{#1}}}
\newcommand{\ImportTok}[1]{#1}
\newcommand{\InformationTok}[1]{\textcolor[rgb]{0.56,0.35,0.01}{\textbf{\textit{#1}}}}
\newcommand{\KeywordTok}[1]{\textcolor[rgb]{0.13,0.29,0.53}{\textbf{#1}}}
\newcommand{\NormalTok}[1]{#1}
\newcommand{\OperatorTok}[1]{\textcolor[rgb]{0.81,0.36,0.00}{\textbf{#1}}}
\newcommand{\OtherTok}[1]{\textcolor[rgb]{0.56,0.35,0.01}{#1}}
\newcommand{\PreprocessorTok}[1]{\textcolor[rgb]{0.56,0.35,0.01}{\textit{#1}}}
\newcommand{\RegionMarkerTok}[1]{#1}
\newcommand{\SpecialCharTok}[1]{\textcolor[rgb]{0.81,0.36,0.00}{\textbf{#1}}}
\newcommand{\SpecialStringTok}[1]{\textcolor[rgb]{0.31,0.60,0.02}{#1}}
\newcommand{\StringTok}[1]{\textcolor[rgb]{0.31,0.60,0.02}{#1}}
\newcommand{\VariableTok}[1]{\textcolor[rgb]{0.00,0.00,0.00}{#1}}
\newcommand{\VerbatimStringTok}[1]{\textcolor[rgb]{0.31,0.60,0.02}{#1}}
\newcommand{\WarningTok}[1]{\textcolor[rgb]{0.56,0.35,0.01}{\textbf{\textit{#1}}}}
\usepackage{longtable,booktabs,array}
\usepackage{calc} % for calculating minipage widths
% Correct order of tables after \paragraph or \subparagraph
\usepackage{etoolbox}
\makeatletter
\patchcmd\longtable{\par}{\if@noskipsec\mbox{}\fi\par}{}{}
\makeatother
% Allow footnotes in longtable head/foot
\IfFileExists{footnotehyper.sty}{\usepackage{footnotehyper}}{\usepackage{footnote}}
\makesavenoteenv{longtable}
\usepackage{graphicx}
\makeatletter
\newsavebox\pandoc@box
\newcommand*\pandocbounded[1]{% scales image to fit in text height/width
  \sbox\pandoc@box{#1}%
  \Gscale@div\@tempa{\textheight}{\dimexpr\ht\pandoc@box+\dp\pandoc@box\relax}%
  \Gscale@div\@tempb{\linewidth}{\wd\pandoc@box}%
  \ifdim\@tempb\p@<\@tempa\p@\let\@tempa\@tempb\fi% select the smaller of both
  \ifdim\@tempa\p@<\p@\scalebox{\@tempa}{\usebox\pandoc@box}%
  \else\usebox{\pandoc@box}%
  \fi%
}
% Set default figure placement to htbp
\def\fps@figure{htbp}
\makeatother
\setlength{\emergencystretch}{3em} % prevent overfull lines
\providecommand{\tightlist}{%
  \setlength{\itemsep}{0pt}\setlength{\parskip}{0pt}}
\usepackage{bookmark}
\IfFileExists{xurl.sty}{\usepackage{xurl}}{} % add URL line breaks if available
\urlstyle{same}
\hypersetup{
  pdftitle={Project 3 Report},
  pdfauthor={Armand Spencer},
  hidelinks,
  pdfcreator={LaTeX via pandoc}}

\title{Project 3 Report}
\usepackage{etoolbox}
\makeatletter
\providecommand{\subtitle}[1]{% add subtitle to \maketitle
  \apptocmd{\@title}{\par {\large #1 \par}}{}{}
}
\makeatother
\subtitle{STATS 220 Semester One 2026}
\author{Armand Spencer}
\date{2026-04-28}

\begin{document}
\maketitle

\subsection{Introduction}\label{introduction}

``Cool'' ``Meerkats''

I picked these two words because I really like meerkats and wanted a
random adjective. So I picked ``Cool''. I felt like if I used a more
niche or distinct adjective then I wouldn't get many good results. Also
meerkats are cool so it's a redundant term in my books.

\href{https://www.pexels.com/search/cool\%20meerkats/}{Cool Meerkats}
\pandocbounded{\includegraphics[keepaspectratio]{screenshotofmeerkats.png}}

All photos generally show meerkats in their natural habitat. Some, but
not many, of the photos show them in enclosures. They are mostly candid
photos, though sometimes the lemur would notice the photographer and
look into the camera. There is also a good mix of photos that are taken
landscape vs portrait. They are also all VERY VERY CUTE.

\begin{Shaded}
\begin{Highlighting}[]
\NormalTok{table }\OtherTok{\textless{}{-}}\NormalTok{ selected\_photos }\SpecialCharTok{\%\textgreater{}\%} \FunctionTok{select}\NormalTok{(url)}
\NormalTok{knitr}\SpecialCharTok{::}\FunctionTok{kable}\NormalTok{(table) }
\end{Highlighting}
\end{Shaded}

\begin{longtable}[]{@{}
  >{\raggedright\arraybackslash}p{(\linewidth - 0\tabcolsep) * \real{1.0000}}@{}}
\toprule\noalign{}
\begin{minipage}[b]{\linewidth}\raggedright
url
\end{minipage} \\
\midrule\noalign{}
\endhead
\bottomrule\noalign{}
\endlastfoot
\url{https://www.pexels.com/photo/close-up-photograph-of-meerkats-11922554/} \\
\url{https://www.pexels.com/photo/meerkats-sitting-on-tree-trunk-17526413/} \\
\url{https://www.pexels.com/photo/meerkats-standing-under-heat-lamp-in-a-warm-setting-29284842/} \\
\url{https://www.pexels.com/photo/family-of-meerkats-in-natural-habitat-37094402/} \\
\url{https://www.pexels.com/photo/a-furry-meerkats-together-7188933/} \\
\url{https://www.pexels.com/photo/close-up-of-meerkats-6272537/} \\
\url{https://www.pexels.com/photo/meerkats-on-rocks-in-sao-paulo-zoo-33640434/} \\
\url{https://www.pexels.com/photo/meerkats-in-the-desert-11025911/} \\
\url{https://www.pexels.com/photo/a-mob-of-meerkats-9450969/} \\
\url{https://www.pexels.com/photo/curious-meerkats-observing-their-surroundings-36484391/} \\
\url{https://www.pexels.com/photo/meerkats-on-a-desert-18157872/} \\
\url{https://www.pexels.com/photo/adorable-meerkats-standing-in-natural-habitat-28656506/} \\
\url{https://www.pexels.com/photo/brown-meerkats-in-close-up-photography-6727776/} \\
\url{https://www.pexels.com/photo/two-meerkats-observing-in-a-natural-habitat-35155650/} \\
\url{https://www.pexels.com/photo/meerkats-beside-a-tree-7722785/} \\
\url{https://www.pexels.com/photo/meerkats-on-the-sand-6796608/} \\
\end{longtable}

\subsection{Key features of my selected
Photos}\label{key-features-of-my-selected-photos}

\begin{Shaded}
\begin{Highlighting}[]
\NormalTok{number\_of\_rows }\OtherTok{\textless{}{-}} \FunctionTok{nrow}\NormalTok{(selected\_photos)}

\NormalTok{large\_image }\OtherTok{\textless{}{-}}\NormalTok{ selected\_photos }\SpecialCharTok{\%\textgreater{}\%}
  \FunctionTok{select}\NormalTok{(width, height) }\SpecialCharTok{\%\textgreater{}\%}
  \FunctionTok{filter}\NormalTok{(width }\SpecialCharTok{\textgreater{}=} \DecValTok{4000} \SpecialCharTok{|}\NormalTok{ height }\SpecialCharTok{\textgreater{}=} \DecValTok{6000}\NormalTok{) }\SpecialCharTok{\%\textgreater{}\%}
  \FunctionTok{nrow}\NormalTok{()}

\NormalTok{photographer\_count }\OtherTok{\textless{}{-}}\NormalTok{ selected\_photos }\SpecialCharTok{\%\textgreater{}\%}
  \FunctionTok{count}\NormalTok{(photographer, }\AttributeTok{sort =} \ConstantTok{TRUE}\NormalTok{)}

\NormalTok{longer\_than\_my\_name }\OtherTok{\textless{}{-}}\NormalTok{ selected\_photos }\SpecialCharTok{\%\textgreater{}\%}
  \FunctionTok{filter}\NormalTok{(}\FunctionTok{str\_detect}\NormalTok{(photographer\_longer\_than\_my\_name, }\StringTok{"nah YEA"}\NormalTok{)) }\SpecialCharTok{\%\textgreater{}\%}
  \FunctionTok{nrow}\NormalTok{()}

\NormalTok{total\_width }\OtherTok{\textless{}{-}}\NormalTok{ selected\_photos}\SpecialCharTok{$}\NormalTok{width }\SpecialCharTok{\%\textgreater{}\%}
  \FunctionTok{sum}\NormalTok{()}
\end{Highlighting}
\end{Shaded}

There are 16 in the selected photo csv. This is due to only 16
photographers using the term ``Meerkats'' in their url.

A Large image is defined as often exceeding 4000x6000 pixels. Of my 16
images, 7 of them fit this description as being a large image for
high-resolution photography

Out of the 16 photographers, the photographer who appeared the most in
my selected\_photos csv was Wayne Jackson, who appeared a total of 2
times.

There are 5 number of photographers whose name is longer then mine is.

Though useless, the sum of all the widths in my data is
\ensuremath{6.8199\times 10^{4}}.

\subsection{Creativity}\label{creativity}

\begin{figure}
\centering
\pandocbounded{\includegraphics[keepaspectratio]{creativity.gif}}
\caption{My Meme}
\end{figure}

I wasn't too creative this time. But in journey to create a meme, I
needed to be able to zoom into the image because I needed to emphasies
that the meerkats is looking in the direction of the camera. So I looked
for another magick function that could let me do this. I really enjoyed
making this meme like usual.

I did have a bit of fun with the ifelse function. Not exactly creative,
but I wrote different strings for if the result is TRUE or FALSE. Not
sure if that counts as creative, but I did that.

Now that I think about it, something different I did do was made a new
variable (via mutate) if the photographers name had more characters than
my name, which is not CRAZY creative but its not common. I also made a
new variable if the photographer calls them ``Meerkats'' or ``Meercats''
and then later only included the rows where they were called Meerkats.

I also had to be a little creative when making the summary values.
Especially if I wanted to find summary values that give somewhat useful
information regarding the dataset. Like I made a new column for if the
photographers names were larger then mine. And then in exploration.R I
had found the mean height grouped by whether their names were bigger
then mine or not.

\subsection{Learning reflection}\label{learning-reflection}

I really like using the ifelse function. One important idea I learnt is
how to manipulate variables in a dataframe in R. I've done something
similar in python and I can say, doing it in R is a lot better and a lot
easier. One thing I also have learnt is how valuable the pipe is, it
makes everything so much easier.

I want to work with databases, in 3A, we briefly touched on SQL and SQL
functions. I would like to expand my knowledge in that region, or
perhaps do some work with relatively humongous databases. Otherwise,
this API stuff I have never touched upon in my life and I didn't know it
exists so I would be happy to take a deeper dive into using APIs.

\subsection{Appendix}\label{appendix}

\begin{Shaded}
\begin{Highlighting}[]
\FunctionTok{library}\NormalTok{(tidyverse)}
\FunctionTok{library}\NormalTok{(httr)}
\FunctionTok{library}\NormalTok{(magick)}

\NormalTok{api\_key }\OtherTok{\textless{}{-}} \StringTok{"w3fvkUmcE3AzHRStvJMJ7ra0zfxyJ97qDpf9rTw5qwLSryDsNB5mk5Jt"}

\NormalTok{url }\OtherTok{\textless{}{-}} \StringTok{"https://api.pexels.com/v1/search?query=cool\%20meerkats\&per\_page=80\&limit=80"}

\NormalTok{response }\OtherTok{\textless{}{-}}\NormalTok{ httr}\SpecialCharTok{::}\FunctionTok{GET}\NormalTok{(url, }
                      \FunctionTok{add\_headers}\NormalTok{(}\AttributeTok{Authorization =}\NormalTok{ api\_key))}

\NormalTok{data }\OtherTok{\textless{}{-}}\NormalTok{ httr}\SpecialCharTok{::}\FunctionTok{content}\NormalTok{(response, }
                      \AttributeTok{as =} \StringTok{"parsed"}\NormalTok{, }
                      \AttributeTok{type =} \StringTok{"application/json"}\NormalTok{)}

\NormalTok{photo\_data }\OtherTok{\textless{}{-}} \FunctionTok{tibble}\NormalTok{(}\AttributeTok{photos =}\NormalTok{ data}\SpecialCharTok{$}\NormalTok{photos) }\SpecialCharTok{\%\textgreater{}\%}
  \FunctionTok{unnest\_wider}\NormalTok{(photos) }\SpecialCharTok{\%\textgreater{}\%}
  \FunctionTok{unnest\_wider}\NormalTok{(src)}

\NormalTok{selected\_photos }\OtherTok{\textless{}{-}}\NormalTok{ photo\_data }\SpecialCharTok{\%\textgreater{}\%} 
  \FunctionTok{mutate}\NormalTok{(}\AttributeTok{photographer\_longer\_than\_my\_name =} \FunctionTok{ifelse}\NormalTok{(}\FunctionTok{str\_count}\NormalTok{(photographer) }\SpecialCharTok{\textgreater{}} \DecValTok{13}\NormalTok{, }\StringTok{"nah YEA"}\NormalTok{, }\StringTok{"Yea NAH nah yea nah"}\NormalTok{)) }\SpecialCharTok{\%\textgreater{}\%} 
  \CommentTok{\# just wanted to see if the photographers name is longer then my name "Armand Spencer" (including spaces)}
  \FunctionTok{mutate}\NormalTok{(}\AttributeTok{name\_words =} \FunctionTok{str\_count}\NormalTok{(photographer, }\StringTok{" "}\NormalTok{) }\SpecialCharTok{+} \DecValTok{1}\NormalTok{) }\SpecialCharTok{\%\textgreater{}\%} 
  \CommentTok{\# counts number of words in photographers name}
  \FunctionTok{mutate}\NormalTok{(}\AttributeTok{meerkats\_or\_meercats =} \FunctionTok{ifelse}\NormalTok{(}\FunctionTok{str\_detect}\NormalTok{(url, }\StringTok{"meerkats"}\NormalTok{), }\StringTok{"MEERKATS"}\NormalTok{, }\StringTok{"weirdos"}\NormalTok{)) }\SpecialCharTok{\%\textgreater{}\%} 
  \CommentTok{\# sees what spelling they use for meerkats}
  \FunctionTok{filter}\NormalTok{(}\FunctionTok{str\_detect}\NormalTok{(meerkats\_or\_meercats, }\StringTok{"MEERKATS"}\NormalTok{))}
 \CommentTok{\# filters based on whether they used MEERKATS or not, cause I only trust people who write Meerkats and NOT Meercats. It looks weird!}
 \CommentTok{\# additionally, out of 80 original rows, only 16 uses meerkats so the 3rd mutate category does do something}
\FunctionTok{write\_csv}\NormalTok{(selected\_photos, }\StringTok{"selected\_photos.csv"}\NormalTok{)}

\CommentTok{\# Summary Values}

\NormalTok{mean\_height }\OtherTok{\textless{}{-}}\NormalTok{ selected\_photos}\SpecialCharTok{$}\NormalTok{height }\SpecialCharTok{\%\textgreater{}\%} \FunctionTok{mean}\NormalTok{(}\AttributeTok{na.rm =} \ConstantTok{TRUE}\NormalTok{) }\SpecialCharTok{\%\textgreater{}\%} \FunctionTok{signif}\NormalTok{(}\DecValTok{4}\NormalTok{) }\CommentTok{\#mean height}
\NormalTok{mean\_width }\OtherTok{\textless{}{-}}\NormalTok{ selected\_photos}\SpecialCharTok{$}\NormalTok{width }\SpecialCharTok{\%\textgreater{}\%} \FunctionTok{mean}\NormalTok{(}\AttributeTok{na.rm =} \ConstantTok{TRUE}\NormalTok{) }\SpecialCharTok{\%\textgreater{}\%} \FunctionTok{signif}\NormalTok{(}\DecValTok{4}\NormalTok{) }\CommentTok{\# mean width}
\NormalTok{mean\_size }\OtherTok{\textless{}{-}} \FunctionTok{paste0}\NormalTok{(}\StringTok{"The mean size of photos in the selected photos are: "}\NormalTok{, mean\_width, }\StringTok{"x"}\NormalTok{, mean\_height) }\CommentTok{\# just shows the mean values together}
\NormalTok{mean\_size}

\NormalTok{max\_name\_size }\OtherTok{\textless{}{-}} \FunctionTok{max}\NormalTok{(selected\_photos}\SpecialCharTok{$}\NormalTok{name\_words) }\CommentTok{\# shows the highest number of words in a name}
\FunctionTok{paste0}\NormalTok{(}\StringTok{"The name with the most number of words has "}\NormalTok{ ,max\_name\_size, }\StringTok{" words in it."}\NormalTok{)}

\NormalTok{summary\_for\_longer\_than\_my\_name }\OtherTok{\textless{}{-}}\NormalTok{ selected\_photos }\SpecialCharTok{\%\textgreater{}\%} 
  \FunctionTok{group\_by}\NormalTok{(photographer\_longer\_than\_my\_name) }\SpecialCharTok{\%\textgreater{}\%} \CommentTok{\#grouped by Yea NAH nah yea nah or nah Yea}
  \FunctionTok{summarise}\NormalTok{(}\AttributeTok{mean\_height\_summary =} \FunctionTok{signif}\NormalTok{(}\FunctionTok{mean}\NormalTok{(height,}\AttributeTok{na.rm =} \ConstantTok{TRUE}\NormalTok{), }\DecValTok{4}\NormalTok{) ) }\CommentTok{\# grabbed mean height of the groups based on Yea NAH or nah Yea}
\NormalTok{nah }\OtherTok{\textless{}{-}}\NormalTok{summary\_for\_longer\_than\_my\_name }\SpecialCharTok{\%\textgreater{}\%} \FunctionTok{select}\NormalTok{(mean\_height\_summary) }\SpecialCharTok{\%\textgreater{}\%} \FunctionTok{slice}\NormalTok{(}\DecValTok{1}\NormalTok{)}
\NormalTok{yea }\OtherTok{\textless{}{-}}\NormalTok{summary\_for\_longer\_than\_my\_name }\SpecialCharTok{\%\textgreater{}\%} \FunctionTok{select}\NormalTok{(mean\_height\_summary) }\SpecialCharTok{\%\textgreater{}\%} \FunctionTok{slice}\NormalTok{(}\DecValTok{2}\NormalTok{)}
\FunctionTok{paste}\NormalTok{(}\StringTok{"The mean height of photos taken by photographers with names longer then mine is:"}\NormalTok{, yea, }\StringTok{"and the mean height of photos taken by photographers }
\StringTok{      with names shorter than mine is:"}\NormalTok{ ,nah)}

\NormalTok{average\_height\_for\_photographer\_shorter\_than\_my\_name }\OtherTok{\textless{}{-}}\NormalTok{ summary\_for\_longer\_than\_my\_name}\SpecialCharTok{$}\NormalTok{mean\_height\_summary[}\DecValTok{1}\NormalTok{] }
\CommentTok{\# grabs the first row, which is Yea Nah nah yea nah, so the average height for names of photographers shorter then mine}
\NormalTok{average\_height\_for\_photographer\_shorter\_than\_my\_name}
\FunctionTok{paste0}\NormalTok{(}\StringTok{"the average height for photographers with names shorter than mine is: "}\NormalTok{, average\_height\_for\_photographer\_shorter\_than\_my\_name)}

\CommentTok{\# Meme MAKING TIME}

\NormalTok{main\_image\_url }\OtherTok{\textless{}{-}}\NormalTok{ selected\_photos}\SpecialCharTok{$}\NormalTok{landscape[}\DecValTok{8}\NormalTok{]}
\NormalTok{main\_image }\OtherTok{\textless{}{-}}\FunctionTok{image\_read}\NormalTok{(}\AttributeTok{path =}\NormalTok{ large\_url)}
\NormalTok{main\_image\_nearly }\OtherTok{\textless{}{-}} \FunctionTok{image\_scale}\NormalTok{(main\_image, }\StringTok{"1200"}\NormalTok{) }
\NormalTok{main\_image\_final }\OtherTok{\textless{}{-}}\NormalTok{ main\_image\_nearly }\SpecialCharTok{\%\textgreater{}\%}
  \FunctionTok{image\_crop}\NormalTok{(}\StringTok{"500x281"}\NormalTok{, }\AttributeTok{gravity =} \StringTok{"center"}\NormalTok{) }\SpecialCharTok{\%\textgreater{}\%} \CommentTok{\# zooms into the lemur face}
  \FunctionTok{image\_scale}\NormalTok{(}\StringTok{"1200"}\NormalTok{) }\CommentTok{\# resizes back to fit meme size}



\NormalTok{secondary\_image\_url }\OtherTok{\textless{}{-}}\NormalTok{ selected\_photos}\SpecialCharTok{$}\NormalTok{landscape[}\DecValTok{10}\NormalTok{]}
\NormalTok{secondary\_image }\OtherTok{\textless{}{-}} \FunctionTok{image\_read}\NormalTok{(}\AttributeTok{path=}\NormalTok{secondary\_image\_url)}
\NormalTok{secondary\_image\_nearly }\OtherTok{\textless{}{-}} \FunctionTok{image\_scale}\NormalTok{(secondary\_image, }\StringTok{"1200"}\NormalTok{)}
\NormalTok{secondary\_image\_final }\OtherTok{\textless{}{-}}\NormalTok{ secondary\_image\_nearly }\SpecialCharTok{\%\textgreater{}\%}
  \FunctionTok{image\_crop}\NormalTok{(}\StringTok{"800x450 + ({-}100) + 200"}\NormalTok{) }\SpecialCharTok{\%\textgreater{}\%} \CommentTok{\# zooms into lemur face again for second image}
  \FunctionTok{image\_scale}\NormalTok{(}\StringTok{"1200"}\NormalTok{) }\CommentTok{\# resizes to fit meme size}
\NormalTok{secondary\_image\_final}

\NormalTok{tertiary\_image\_url }\OtherTok{\textless{}{-}}\NormalTok{ selected\_photos}\SpecialCharTok{$}\NormalTok{landscape[}\DecValTok{12}\NormalTok{]}
\NormalTok{tertiary\_image }\OtherTok{\textless{}{-}} \FunctionTok{image\_read}\NormalTok{(}\AttributeTok{path=}\NormalTok{tertiary\_image\_url)}
\NormalTok{tertiary\_image\_final }\OtherTok{\textless{}{-}} \FunctionTok{image\_scale}\NormalTok{(tertiary\_image, }\StringTok{"1200"}\NormalTok{)}


\NormalTok{text1 }\OtherTok{\textless{}{-}} \FunctionTok{image\_blank}\NormalTok{(}\AttributeTok{width =}\DecValTok{1200}\NormalTok{,}
                     \AttributeTok{height =}\DecValTok{200}\NormalTok{, }
                     \AttributeTok{color =} \StringTok{"white"}\NormalTok{) }\SpecialCharTok{\%\textgreater{}\%}
  \FunctionTok{image\_annotate}\NormalTok{(}\AttributeTok{text =} \StringTok{"Me : Yo bro, DON\textquotesingle{}T turn around but{-}"}\NormalTok{,}
                 \AttributeTok{gravity =} \StringTok{"Center"}\NormalTok{, }
                 \AttributeTok{size =} \DecValTok{40}\NormalTok{,}
                 \AttributeTok{color =} \StringTok{"black"}\NormalTok{,}
                 \AttributeTok{font =} \StringTok{"Impact"}\NormalTok{)}

\NormalTok{text2 }\OtherTok{\textless{}{-}} \FunctionTok{image\_blank}\NormalTok{(}\AttributeTok{width =}\DecValTok{1200}\NormalTok{,}
                     \AttributeTok{height =}\DecValTok{200}\NormalTok{, }
                     \AttributeTok{color =} \StringTok{"white"}\NormalTok{) }\SpecialCharTok{\%\textgreater{}\%}
  \FunctionTok{image\_annotate}\NormalTok{(}\AttributeTok{text =} \StringTok{"Bro 0.00001s later"}\NormalTok{,}
                 \AttributeTok{gravity =} \StringTok{"Center"}\NormalTok{, }
                 \AttributeTok{size =} \DecValTok{40}\NormalTok{,}
                 \AttributeTok{color =} \StringTok{"black"}\NormalTok{,}
                 \AttributeTok{font =} \StringTok{"Impact"}\NormalTok{)}

\NormalTok{frame1 }\OtherTok{\textless{}{-}} \FunctionTok{c}\NormalTok{(text1, tertiary\_image\_final)}
\NormalTok{frame2 }\OtherTok{\textless{}{-}} \FunctionTok{c}\NormalTok{(text2, secondary\_image\_nearly)}
\NormalTok{frame3 }\OtherTok{\textless{}{-}} \FunctionTok{c}\NormalTok{(text2, secondary\_image\_final)}
\NormalTok{frame4 }\OtherTok{\textless{}{-}} \FunctionTok{c}\NormalTok{(text2, main\_image\_nearly)}
\NormalTok{frame5 }\OtherTok{\textless{}{-}} \FunctionTok{c}\NormalTok{(text2, main\_image\_final) }\CommentTok{\# creates the image vector}

\NormalTok{frame1\_f }\OtherTok{\textless{}{-}} \FunctionTok{image\_append}\NormalTok{(frame1, }\AttributeTok{stack =} \ConstantTok{TRUE}\NormalTok{)}
\NormalTok{frame2\_f }\OtherTok{\textless{}{-}} \FunctionTok{image\_append}\NormalTok{(frame2, }\AttributeTok{stack =} \ConstantTok{TRUE}\NormalTok{)}
\NormalTok{frame3\_f }\OtherTok{\textless{}{-}} \FunctionTok{image\_append}\NormalTok{(frame3, }\AttributeTok{stack =} \ConstantTok{TRUE}\NormalTok{)}
\NormalTok{frame4\_f }\OtherTok{\textless{}{-}} \FunctionTok{image\_append}\NormalTok{(frame4, }\AttributeTok{stack =} \ConstantTok{TRUE}\NormalTok{)}
\NormalTok{frame5\_f }\OtherTok{\textless{}{-}} \FunctionTok{image\_append}\NormalTok{(frame5, }\AttributeTok{stack =} \ConstantTok{TRUE}\NormalTok{) }\CommentTok{\# adds the two images together}

\NormalTok{animation }\OtherTok{\textless{}{-}} \FunctionTok{c}\NormalTok{(frame1\_f, frame2\_f, frame3\_f, frame4\_f, frame5\_f)}

\NormalTok{lemur\_meme }\OtherTok{\textless{}{-}} \FunctionTok{image\_animate}\NormalTok{(animation, }\AttributeTok{fps=}\DecValTok{1}\NormalTok{) }\CommentTok{\# gif completed}

\NormalTok{upload }\OtherTok{\textless{}{-}}\FunctionTok{image\_write}\NormalTok{(lemur\_meme, }\StringTok{"creativity.gif"}\NormalTok{)}
\end{Highlighting}
\end{Shaded}

\begin{Shaded}
\begin{Highlighting}[]
\SpecialCharTok{{-}{-}{-}}
\NormalTok{title}\SpecialCharTok{:} \StringTok{"Project 3 Report"}
\NormalTok{author}\SpecialCharTok{:} \StringTok{"Armand Spencer"}
\NormalTok{date}\SpecialCharTok{:} \StringTok{"2026{-}04{-}28"}
\NormalTok{output}\SpecialCharTok{:}
\NormalTok{  pdf\_document}\SpecialCharTok{:}\NormalTok{ default}
\NormalTok{  html\_document}\SpecialCharTok{:}
\NormalTok{    code folding}\SpecialCharTok{:}\NormalTok{ hide}
\NormalTok{subtitle}\SpecialCharTok{:}\NormalTok{ STATS }\DecValTok{220}\NormalTok{ Semester One }\DecValTok{2026}
\SpecialCharTok{{-}{-}{-}}

\StringTok{\textasciigrave{}\textasciigrave{}\textasciigrave{}}\AttributeTok{\{r setup, include=FALSE\}}
\AttributeTok{knitr::opts\_chunk$set(echo=TRUE, message=FALSE, warning=FALSE, error=FALSE)}
\AttributeTok{library(tidyverse)}
\AttributeTok{library(knitr)}

\AttributeTok{selected\_photos \textless{}{-} read\_csv("selected\_photos.csv")}
\StringTok{\textasciigrave{}\textasciigrave{}\textasciigrave{}}

\StringTok{\textasciigrave{}\textasciigrave{}\textasciigrave{}}\AttributeTok{\{css echo=FALSE\}}

\AttributeTok{body \{}
\AttributeTok{  font{-}family:  "Times New Roman", serif;}
\AttributeTok{  line{-}height: 1.6;}
\AttributeTok{  max{-}width: 900px;}
\AttributeTok{  margin: auto;}
\AttributeTok{  padding: 20px;}
\AttributeTok{  background: linear{-}gradient(155deg, \#B5E2FA, \#F9F7F3,\#EDDEA4, \#F7A072 );}
\AttributeTok{  font{-}size: 15px;}
\AttributeTok{\}}
\AttributeTok{p \{}
\AttributeTok{    font{-}family: "Times New Roman", Helvetica, Roboto, sans{-}serif;}
\AttributeTok{    font{-}size: 15px;}
\AttributeTok{\}}

\AttributeTok{h1 \{}
\AttributeTok{  font{-}family: Inter,  Helvetica , Roboto, sans{-}serif;}
\AttributeTok{  color: \#2C3E50;}
\AttributeTok{  font{-}weight: 600;}
\AttributeTok{\}}

\AttributeTok{h2 \{}
\AttributeTok{  font{-}family: Inter, Helvetica , Roboto, sans{-}serif;}
\AttributeTok{  color: \#2C3E50;}
\AttributeTok{  font{-}weight: 600;}
\AttributeTok{\}}

\AttributeTok{h3 \{}
\AttributeTok{  font{-}family:  Inter, Helvetica, Roboto, sans{-}serif;}
\AttributeTok{  color: \#2C3E50;}
\AttributeTok{  font{-}weight: 600;}
\AttributeTok{\}}

\AttributeTok{img \{}
\AttributeTok{  max{-}width: 100\%;}
\AttributeTok{  height: auto;}
\AttributeTok{  display: block;}
\AttributeTok{  margin: 0 auto;}
\AttributeTok{  }
\AttributeTok{\}}

\StringTok{\textasciigrave{}\textasciigrave{}\textasciigrave{}}
\DocumentationTok{\#\# Introduction}

\StringTok{"Cool"} \StringTok{"Meerkats"}

\NormalTok{I picked these two words because I really like meerkats and wanted a random adjective. So I picked }\StringTok{"Cool"}\NormalTok{. I felt like }\ControlFlowTok{if}\NormalTok{ I used a more niche or distinct adjective then I wouldn}\StringTok{\textquotesingle{}t get many good results. Also meerkats are cool so it\textquotesingle{}}\NormalTok{s a redundant term }\ControlFlowTok{in}\NormalTok{ my books.}

\NormalTok{[Cool Meerkats](https}\SpecialCharTok{:}\ErrorTok{//}\NormalTok{www.pexels.com}\SpecialCharTok{/}\NormalTok{search}\SpecialCharTok{/}\NormalTok{cool\%}\DecValTok{20}\NormalTok{meerkats}\SpecialCharTok{/}\NormalTok{) }
\SpecialCharTok{!}\NormalTok{[Image of Meerkats](screenshotofmeerkats.png)}

\NormalTok{All photos generally show meerkats }\ControlFlowTok{in}\NormalTok{ their natural habitat. Some, but not many, of the photos show them }\ControlFlowTok{in}\NormalTok{ enclosures. They are mostly candid photos, though sometimes the lemur would notice the photographer and look into the camera. There is also a good mix of photos that are taken landscape vs portrait. They are also all VERY VERY CUTE.}

\StringTok{\textasciigrave{}\textasciigrave{}\textasciigrave{}}\AttributeTok{\{r\}}
\AttributeTok{table \textless{}{-} selected\_photos \%\textgreater{}\% select(url)}
\AttributeTok{knitr::kable(table) }
\StringTok{\textasciigrave{}\textasciigrave{}\textasciigrave{}}

\DocumentationTok{\#\# Key features of my selected Photos}
\StringTok{\textasciigrave{}\textasciigrave{}\textasciigrave{}}\AttributeTok{\{r\}}
\AttributeTok{number\_of\_rows \textless{}{-} nrow(selected\_photos)}

\AttributeTok{large\_image \textless{}{-} selected\_photos \%\textgreater{}\%}
\AttributeTok{  select(width, height) \%\textgreater{}\%}
\AttributeTok{  filter(width \textgreater{}= 4000 | height \textgreater{}= 6000) \%\textgreater{}\%}
\AttributeTok{  nrow()}

\AttributeTok{photographer\_count \textless{}{-} selected\_photos \%\textgreater{}\%}
\AttributeTok{  count(photographer, sort = TRUE)}

\AttributeTok{longer\_than\_my\_name \textless{}{-} selected\_photos \%\textgreater{}\%}
\AttributeTok{  filter(str\_detect(photographer\_longer\_than\_my\_name, "nah YEA")) \%\textgreater{}\%}
\AttributeTok{  nrow()}

\AttributeTok{total\_width \textless{}{-} selected\_photos$width \%\textgreater{}\%}
\AttributeTok{  sum()}

\StringTok{\textasciigrave{}\textasciigrave{}\textasciigrave{}}
\NormalTok{There are }\StringTok{\textasciigrave{}}\AttributeTok{r number\_of\_rows}\StringTok{\textasciigrave{}} \ControlFlowTok{in}\NormalTok{ the selected photo csv. This is due to only }\DecValTok{16}\NormalTok{ photographers using the term }\StringTok{"Meerkats"} \ControlFlowTok{in}\NormalTok{ their url.}

\NormalTok{A Large image is defined as often exceeding }\DecValTok{4000}\NormalTok{x6000 pixels. Of my }\DecValTok{16}\NormalTok{ images, }\StringTok{\textasciigrave{}}\AttributeTok{r large\_image}\StringTok{\textasciigrave{}}\NormalTok{ of them fit this description as being a large image }\ControlFlowTok{for}\NormalTok{ high}\SpecialCharTok{{-}}\NormalTok{resolution photography}

\NormalTok{Out of the }\DecValTok{16}\NormalTok{ photographers, the photographer who appeared the most }\ControlFlowTok{in}\NormalTok{ my selected\_photos csv was }\StringTok{\textasciigrave{}}\AttributeTok{r photographer\_count \%\textgreater{}\% select(photographer) \%\textgreater{}\% slice(1)}\StringTok{\textasciigrave{}}\NormalTok{, who appeared a total of }\StringTok{\textasciigrave{}}\AttributeTok{r photographer\_count \%\textgreater{}\% select(n) \%\textgreater{}\% slice(1)}\StringTok{\textasciigrave{}}\NormalTok{ times.}

\NormalTok{There are }\StringTok{\textasciigrave{}}\AttributeTok{r longer\_than\_my\_name}\StringTok{\textasciigrave{}}\NormalTok{ number of photographers whose name is longer then mine is.}

\NormalTok{Though useless, the sum of all the widths }\ControlFlowTok{in}\NormalTok{ my data is }\StringTok{\textasciigrave{}}\AttributeTok{r total\_width}\StringTok{\textasciigrave{}}\NormalTok{.}

\DocumentationTok{\#\# Creativity}

\SpecialCharTok{!}\NormalTok{[My Meme](creativity.gif)}

\NormalTok{I wasn}\StringTok{\textquotesingle{}t too creative this time. But in journey to create a meme, I needed to be able to zoom into the image because I needed to emphasies that the meerkats is looking in the direction of the camera. So I looked for another magick function that could let me do this. I really enjoyed making this meme like usual.}

\StringTok{I did have a bit of fun with the ifelse function. Not exactly creative, but I wrote different strings for if the result is TRUE or FALSE. Not sure if that counts as creative, but I did that.}

\StringTok{Now that I think about it, something different I did do was made a new variable (via mutate) if the photographers name had more characters than my name, which is not CRAZY creative but its not common. I also made a new variable if the photographer calls them "Meerkats" or "Meercats" and then later only included the rows where they were called Meerkats.}

\StringTok{I also had to be a little creative when making the summary values. Especially if I wanted to find summary values that give somewhat useful information regarding the dataset. Like I made a new column for if the photographers names were larger then mine. And then in exploration.R I had found the mean height grouped by whether their names were bigger then mine or not. }


\StringTok{\#\# Learning reflection}

\StringTok{I really like using the ifelse function. One important idea I learnt is how to manipulate variables in a dataframe in R. I\textquotesingle{}}\NormalTok{ve done something similar }\ControlFlowTok{in}\NormalTok{ python and I can say, doing it }\ControlFlowTok{in}\NormalTok{ R is a lot better and a lot easier. One thing I also have learnt is how valuable the pipe is, it makes everything so much easier.}

\NormalTok{I want to work with databases, }\ControlFlowTok{in} \DecValTok{3}\NormalTok{A, we briefly touched on SQL and SQL functions. I would like to expand my knowledge }\ControlFlowTok{in}\NormalTok{ that region, or perhaps do some work with relatively humongous databases. Otherwise, this API stuff I have never touched upon }\ControlFlowTok{in}\NormalTok{ my life and I didn}\StringTok{\textquotesingle{}t know it exists so I would be happy to take a deeper dive into using APIs.}

\StringTok{\#\# Appendix}

\StringTok{\textasciigrave{}\textasciigrave{}\textasciigrave{}\{r file=\textquotesingle{}}\NormalTok{exploration.R}\StringTok{\textquotesingle{}, eval=FALSE, echo=TRUE\}}

\StringTok{\textasciigrave{}\textasciigrave{}\textasciigrave{}}
\StringTok{\textasciigrave{}\textasciigrave{}\textasciigrave{}\{r file=\textquotesingle{}}\NormalTok{project3\_report.Rmd}\StringTok{\textquotesingle{}, eval=FALSE, echo=TRUE\}}

\StringTok{\textasciigrave{}\textasciigrave{}\textasciigrave{}}
\end{Highlighting}
\end{Shaded}


\end{document}
